% ==============================================================================
% beamer-header.tex -- Beamer preamble for research talks
%
% The house style is plain. Default theme, white background, no colored header
% bar or footline, 9pt, 4:3. Figures carry most of the argument, frames often
% pair an image column with a text column, and emphasis is done with muted text
% color.
%
% Two parameters, each a one-word change in the \documentclass line.
%   aspect ratio  4:3 by default. Add aspectratio=169 for widescreen.
%   font size     9pt by default. Use 10pt or 11pt if a room needs it.
%
% Frame content is vertically centered, the beamer default. For a long course
% deck add the class option t, as in \documentclass[9pt,t]{beamer}, so content
% is top-aligned and stays put from slide to slide.
%
% Usage. A deck in Slides/ begins with % ==============================================================================
% beamer-header.tex -- Beamer preamble for research talks
%
% The house style is plain. Default theme, white background, no colored header
% bar or footline, 9pt, 4:3. Figures carry most of the argument, frames often
% pair an image column with a text column, and emphasis is done with muted text
% color.
%
% Two parameters, each a one-word change in the \documentclass line.
%   aspect ratio  4:3 by default. Add aspectratio=169 for widescreen.
%   font size     9pt by default. Use 10pt or 11pt if a room needs it.
%
% Frame content is vertically centered, the beamer default. For a long course
% deck add the class option t, as in \documentclass[9pt,t]{beamer}, so content
% is top-aligned and stays put from slide to slide.
%
% Usage. A deck in Slides/ begins with % ==============================================================================
% beamer-header.tex -- Beamer preamble for research talks
%
% The house style is plain. Default theme, white background, no colored header
% bar or footline, 9pt, 4:3. Figures carry most of the argument, frames often
% pair an image column with a text column, and emphasis is done with muted text
% color.
%
% Two parameters, each a one-word change in the \documentclass line.
%   aspect ratio  4:3 by default. Add aspectratio=169 for widescreen.
%   font size     9pt by default. Use 10pt or 11pt if a room needs it.
%
% Frame content is vertically centered, the beamer default. For a long course
% deck add the class option t, as in \documentclass[9pt,t]{beamer}, so content
% is top-aligned and stays put from slide to slide.
%
% Usage. A deck in Slides/ begins with % ==============================================================================
% beamer-header.tex -- Beamer preamble for research talks
%
% The house style is plain. Default theme, white background, no colored header
% bar or footline, 9pt, 4:3. Figures carry most of the argument, frames often
% pair an image column with a text column, and emphasis is done with muted text
% color.
%
% Two parameters, each a one-word change in the \documentclass line.
%   aspect ratio  4:3 by default. Add aspectratio=169 for widescreen.
%   font size     9pt by default. Use 10pt or 11pt if a room needs it.
%
% Frame content is vertically centered, the beamer default. For a long course
% deck add the class option t, as in \documentclass[9pt,t]{beamer}, so content
% is top-aligned and stays put from slide to slide.
%
% Usage. A deck in Slides/ begins with \input{../Preambles/beamer-header.tex}
% and compiles from Slides/ with
%   xelatex -interaction=nonstopmode deck.tex
%   BIBINPUTS=..:$BIBINPUTS bibtex deck     (only if the deck cites)
%   xelatex -interaction=nonstopmode deck.tex
% XeLaTeX is required because fontspec is loaded below. Figures resolve from
% Figures/, so \includegraphics{fig1} finds Figures/fig1.pdf.
% ==============================================================================

\documentclass[9pt]{beamer}

% ------------------------------------------------------------------------------
% Theme
% ------------------------------------------------------------------------------
\usetheme{default}
\usecolortheme{default}
\setbeamertemplate{navigation symbols}{}
\setbeamertemplate{itemize items}[circle]

% ------------------------------------------------------------------------------
% Packages
% ------------------------------------------------------------------------------
\usepackage{amsmath,amssymb,tcolorbox}
\usepackage{listings}
\usepackage{multirow}
\usepackage{booktabs}
\usepackage{graphicx}
% adjustbox[export] adds the frame key to \includegraphics, so
% \includegraphics[width=\linewidth, frame]{...} draws a thin border.
\usepackage[export]{adjustbox}
\usepackage{xcolor}
\usepackage{fontspec}
\usepackage[authoryear,round]{natbib}
\usepackage{hyperref}
\hypersetup{
    colorlinks=true,
    linkcolor=blue,
    filecolor=magenta,
    urlcolor=cyan,
}
\usepackage{caption}
% Removes the "Figure:" prefix from captions.
\captionsetup[figure]{name=}

\graphicspath{{../Figures/}{./Figures/}}

% ------------------------------------------------------------------------------
% Palette. Emphasis uses the two muted text colors below. The named colors let
% TikZ diagrams reuse the same palette.
% ------------------------------------------------------------------------------
\definecolor{houseblue}{RGB}{31,58,110}    % matches blue!70!black
\definecolor{housered}{RGB}{140,27,27}     % matches red!70!black

% ------------------------------------------------------------------------------
% Emphasis
%   \term{...}   a defined term or key phrase, bold muted blue
%   \alert{...}  a warning or the one thing to notice, bold muted red
%   \muted{...}  secondary text, gray
%   \code{...}   inline code
% ------------------------------------------------------------------------------
\newcommand{\term}[1]{\textbf{\textcolor{blue!70!black}{#1}}}
\renewcommand{\alert}[1]{\textbf{\textcolor{red!70!black}{#1}}}
\newcommand{\muted}[1]{\textcolor{gray}{#1}}
\newcommand{\code}[1]{\texttt{#1}}

% ------------------------------------------------------------------------------
% Title frame
%   \titleframe{Title}{Authors and affiliations}{Venue and date}{email}
% ------------------------------------------------------------------------------
\newcommand{\titleframe}[4]{%
  \begin{frame}[c]
    \centering
    \huge{#1}\\
    \vspace{0.25in}
    \normalsize{#2}\\
    \vspace{0.25in}
    \footnotesize{#3}\\
    \vspace{.15in}
    \footnotesize{This Version: \today}\\
    \vspace{.25in}
    \href{mailto:#4}{#4}
  \end{frame}%
}

% ------------------------------------------------------------------------------
% Figure frames
% ------------------------------------------------------------------------------

% \source{...}  small gray attribution line under an image
\newcommand{\source}[1]{%
  \par\vspace{2pt}%
  {\scriptsize\textcolor{gray}{Source: #1}\par}%
}

% \figframe{Title}{graphic}{caption}   one figure in a titled frame
\newcommand{\figframe}[3]{%
  \begin{frame}{#1}
    \begin{center}
      \includegraphics[width=\textwidth,height=0.80\textheight,keepaspectratio]{#2}
    \end{center}
    {\scriptsize\textcolor{gray}{#3}\par}
  \end{frame}%
}

% \fullfigframe{graphic}{caption}      one figure, no title, maximum size
\newcommand{\fullfigframe}[2]{%
  \begin{frame}[plain]
    \begin{center}
      \includegraphics[width=\textwidth,height=0.90\textheight,keepaspectratio]{#1}
    \end{center}
    {\scriptsize\textcolor{gray}{#2}\par}
  \end{frame}%
}

% ------------------------------------------------------------------------------
% keybox. One quiet box for the single takeaway of a frame. Use it rarely.
% ------------------------------------------------------------------------------
\newtcolorbox{keybox}{
  colback=black!3, colframe=orange!55!black,
  boxrule=0.6pt, arc=1pt,
  left=6pt, right=6pt, top=4pt, bottom=4pt
}

% ------------------------------------------------------------------------------
% Listings
% ------------------------------------------------------------------------------
\lstset{
  basicstyle=\ttfamily\scriptsize,
  keywordstyle=\color{blue!70!black}\bfseries,
  commentstyle=\color{gray}\itshape,
  stringstyle=\color{green!40!black},
  frame=single,
  framerule=0.4pt,
  rulecolor=\color{gray!50},
  breaklines=true,
  numbers=none,
  showstringspaces=false,
  tabsize=2
}

% ------------------------------------------------------------------------------
% Section dividers. Each \section{...} produces a plain frame with its title.
% ------------------------------------------------------------------------------
\AtBeginSection[]{%
  \begin{frame}[plain,c]
    \begin{center}
      {\usebeamercolor[fg]{structure}\Large\bfseries\insertsection}
    \end{center}
  \end{frame}%
}

% and compiles from Slides/ with
%   xelatex -interaction=nonstopmode deck.tex
%   BIBINPUTS=..:$BIBINPUTS bibtex deck     (only if the deck cites)
%   xelatex -interaction=nonstopmode deck.tex
% XeLaTeX is required because fontspec is loaded below. Figures resolve from
% Figures/, so \includegraphics{fig1} finds Figures/fig1.pdf.
% ==============================================================================

\documentclass[9pt]{beamer}

% ------------------------------------------------------------------------------
% Theme
% ------------------------------------------------------------------------------
\usetheme{default}
\usecolortheme{default}
\setbeamertemplate{navigation symbols}{}
\setbeamertemplate{itemize items}[circle]

% ------------------------------------------------------------------------------
% Packages
% ------------------------------------------------------------------------------
\usepackage{amsmath,amssymb,tcolorbox}
\usepackage{listings}
\usepackage{multirow}
\usepackage{booktabs}
\usepackage{graphicx}
% adjustbox[export] adds the frame key to \includegraphics, so
% \includegraphics[width=\linewidth, frame]{...} draws a thin border.
\usepackage[export]{adjustbox}
\usepackage{xcolor}
\usepackage{fontspec}
\usepackage[authoryear,round]{natbib}
\usepackage{hyperref}
\hypersetup{
    colorlinks=true,
    linkcolor=blue,
    filecolor=magenta,
    urlcolor=cyan,
}
\usepackage{caption}
% Removes the "Figure:" prefix from captions.
\captionsetup[figure]{name=}

\graphicspath{{../Figures/}{./Figures/}}

% ------------------------------------------------------------------------------
% Palette. Emphasis uses the two muted text colors below. The named colors let
% TikZ diagrams reuse the same palette.
% ------------------------------------------------------------------------------
\definecolor{houseblue}{RGB}{31,58,110}    % matches blue!70!black
\definecolor{housered}{RGB}{140,27,27}     % matches red!70!black

% ------------------------------------------------------------------------------
% Emphasis
%   \term{...}   a defined term or key phrase, bold muted blue
%   \alert{...}  a warning or the one thing to notice, bold muted red
%   \muted{...}  secondary text, gray
%   \code{...}   inline code
% ------------------------------------------------------------------------------
\newcommand{\term}[1]{\textbf{\textcolor{blue!70!black}{#1}}}
\renewcommand{\alert}[1]{\textbf{\textcolor{red!70!black}{#1}}}
\newcommand{\muted}[1]{\textcolor{gray}{#1}}
\newcommand{\code}[1]{\texttt{#1}}

% ------------------------------------------------------------------------------
% Title frame
%   \titleframe{Title}{Authors and affiliations}{Venue and date}{email}
% ------------------------------------------------------------------------------
\newcommand{\titleframe}[4]{%
  \begin{frame}[c]
    \centering
    \huge{#1}\\
    \vspace{0.25in}
    \normalsize{#2}\\
    \vspace{0.25in}
    \footnotesize{#3}\\
    \vspace{.15in}
    \footnotesize{This Version: \today}\\
    \vspace{.25in}
    \href{mailto:#4}{#4}
  \end{frame}%
}

% ------------------------------------------------------------------------------
% Figure frames
% ------------------------------------------------------------------------------

% \source{...}  small gray attribution line under an image
\newcommand{\source}[1]{%
  \par\vspace{2pt}%
  {\scriptsize\textcolor{gray}{Source: #1}\par}%
}

% \figframe{Title}{graphic}{caption}   one figure in a titled frame
\newcommand{\figframe}[3]{%
  \begin{frame}{#1}
    \begin{center}
      \includegraphics[width=\textwidth,height=0.80\textheight,keepaspectratio]{#2}
    \end{center}
    {\scriptsize\textcolor{gray}{#3}\par}
  \end{frame}%
}

% \fullfigframe{graphic}{caption}      one figure, no title, maximum size
\newcommand{\fullfigframe}[2]{%
  \begin{frame}[plain]
    \begin{center}
      \includegraphics[width=\textwidth,height=0.90\textheight,keepaspectratio]{#1}
    \end{center}
    {\scriptsize\textcolor{gray}{#2}\par}
  \end{frame}%
}

% ------------------------------------------------------------------------------
% keybox. One quiet box for the single takeaway of a frame. Use it rarely.
% ------------------------------------------------------------------------------
\newtcolorbox{keybox}{
  colback=black!3, colframe=orange!55!black,
  boxrule=0.6pt, arc=1pt,
  left=6pt, right=6pt, top=4pt, bottom=4pt
}

% ------------------------------------------------------------------------------
% Listings
% ------------------------------------------------------------------------------
\lstset{
  basicstyle=\ttfamily\scriptsize,
  keywordstyle=\color{blue!70!black}\bfseries,
  commentstyle=\color{gray}\itshape,
  stringstyle=\color{green!40!black},
  frame=single,
  framerule=0.4pt,
  rulecolor=\color{gray!50},
  breaklines=true,
  numbers=none,
  showstringspaces=false,
  tabsize=2
}

% ------------------------------------------------------------------------------
% Section dividers. Each \section{...} produces a plain frame with its title.
% ------------------------------------------------------------------------------
\AtBeginSection[]{%
  \begin{frame}[plain,c]
    \begin{center}
      {\usebeamercolor[fg]{structure}\Large\bfseries\insertsection}
    \end{center}
  \end{frame}%
}

% and compiles from Slides/ with
%   xelatex -interaction=nonstopmode deck.tex
%   BIBINPUTS=..:$BIBINPUTS bibtex deck     (only if the deck cites)
%   xelatex -interaction=nonstopmode deck.tex
% XeLaTeX is required because fontspec is loaded below. Figures resolve from
% Figures/, so \includegraphics{fig1} finds Figures/fig1.pdf.
% ==============================================================================

\documentclass[9pt]{beamer}

% ------------------------------------------------------------------------------
% Theme
% ------------------------------------------------------------------------------
\usetheme{default}
\usecolortheme{default}
\setbeamertemplate{navigation symbols}{}
\setbeamertemplate{itemize items}[circle]

% ------------------------------------------------------------------------------
% Packages
% ------------------------------------------------------------------------------
\usepackage{amsmath,amssymb,tcolorbox}
\usepackage{listings}
\usepackage{multirow}
\usepackage{booktabs}
\usepackage{graphicx}
% adjustbox[export] adds the frame key to \includegraphics, so
% \includegraphics[width=\linewidth, frame]{...} draws a thin border.
\usepackage[export]{adjustbox}
\usepackage{xcolor}
\usepackage{fontspec}
\usepackage[authoryear,round]{natbib}
\usepackage{hyperref}
\hypersetup{
    colorlinks=true,
    linkcolor=blue,
    filecolor=magenta,
    urlcolor=cyan,
}
\usepackage{caption}
% Removes the "Figure:" prefix from captions.
\captionsetup[figure]{name=}

\graphicspath{{../Figures/}{./Figures/}}

% ------------------------------------------------------------------------------
% Palette. Emphasis uses the two muted text colors below. The named colors let
% TikZ diagrams reuse the same palette.
% ------------------------------------------------------------------------------
\definecolor{houseblue}{RGB}{31,58,110}    % matches blue!70!black
\definecolor{housered}{RGB}{140,27,27}     % matches red!70!black

% ------------------------------------------------------------------------------
% Emphasis
%   \term{...}   a defined term or key phrase, bold muted blue
%   \alert{...}  a warning or the one thing to notice, bold muted red
%   \muted{...}  secondary text, gray
%   \code{...}   inline code
% ------------------------------------------------------------------------------
\newcommand{\term}[1]{\textbf{\textcolor{blue!70!black}{#1}}}
\renewcommand{\alert}[1]{\textbf{\textcolor{red!70!black}{#1}}}
\newcommand{\muted}[1]{\textcolor{gray}{#1}}
\newcommand{\code}[1]{\texttt{#1}}

% ------------------------------------------------------------------------------
% Title frame
%   \titleframe{Title}{Authors and affiliations}{Venue and date}{email}
% ------------------------------------------------------------------------------
\newcommand{\titleframe}[4]{%
  \begin{frame}[c]
    \centering
    \huge{#1}\\
    \vspace{0.25in}
    \normalsize{#2}\\
    \vspace{0.25in}
    \footnotesize{#3}\\
    \vspace{.15in}
    \footnotesize{This Version: \today}\\
    \vspace{.25in}
    \href{mailto:#4}{#4}
  \end{frame}%
}

% ------------------------------------------------------------------------------
% Figure frames
% ------------------------------------------------------------------------------

% \source{...}  small gray attribution line under an image
\newcommand{\source}[1]{%
  \par\vspace{2pt}%
  {\scriptsize\textcolor{gray}{Source: #1}\par}%
}

% \figframe{Title}{graphic}{caption}   one figure in a titled frame
\newcommand{\figframe}[3]{%
  \begin{frame}{#1}
    \begin{center}
      \includegraphics[width=\textwidth,height=0.80\textheight,keepaspectratio]{#2}
    \end{center}
    {\scriptsize\textcolor{gray}{#3}\par}
  \end{frame}%
}

% \fullfigframe{graphic}{caption}      one figure, no title, maximum size
\newcommand{\fullfigframe}[2]{%
  \begin{frame}[plain]
    \begin{center}
      \includegraphics[width=\textwidth,height=0.90\textheight,keepaspectratio]{#1}
    \end{center}
    {\scriptsize\textcolor{gray}{#2}\par}
  \end{frame}%
}

% ------------------------------------------------------------------------------
% keybox. One quiet box for the single takeaway of a frame. Use it rarely.
% ------------------------------------------------------------------------------
\newtcolorbox{keybox}{
  colback=black!3, colframe=orange!55!black,
  boxrule=0.6pt, arc=1pt,
  left=6pt, right=6pt, top=4pt, bottom=4pt
}

% ------------------------------------------------------------------------------
% Listings
% ------------------------------------------------------------------------------
\lstset{
  basicstyle=\ttfamily\scriptsize,
  keywordstyle=\color{blue!70!black}\bfseries,
  commentstyle=\color{gray}\itshape,
  stringstyle=\color{green!40!black},
  frame=single,
  framerule=0.4pt,
  rulecolor=\color{gray!50},
  breaklines=true,
  numbers=none,
  showstringspaces=false,
  tabsize=2
}

% ------------------------------------------------------------------------------
% Section dividers. Each \section{...} produces a plain frame with its title.
% ------------------------------------------------------------------------------
\AtBeginSection[]{%
  \begin{frame}[plain,c]
    \begin{center}
      {\usebeamercolor[fg]{structure}\Large\bfseries\insertsection}
    \end{center}
  \end{frame}%
}

% and compiles from Slides/ with
%   xelatex -interaction=nonstopmode deck.tex
%   BIBINPUTS=..:$BIBINPUTS bibtex deck     (only if the deck cites)
%   xelatex -interaction=nonstopmode deck.tex
% XeLaTeX is required because fontspec is loaded below. Figures resolve from
% Figures/, so \includegraphics{fig1} finds Figures/fig1.pdf.
% ==============================================================================

\documentclass[9pt]{beamer}

% ------------------------------------------------------------------------------
% Theme
% ------------------------------------------------------------------------------
\usetheme{default}
\usecolortheme{default}
\setbeamertemplate{navigation symbols}{}
\setbeamertemplate{itemize items}[circle]

% ------------------------------------------------------------------------------
% Packages
% ------------------------------------------------------------------------------
\usepackage{amsmath,amssymb,tcolorbox}
\usepackage{listings}
\usepackage{multirow}
\usepackage{booktabs}
\usepackage{graphicx}
% adjustbox[export] adds the frame key to \includegraphics, so
% \includegraphics[width=\linewidth, frame]{...} draws a thin border.
\usepackage[export]{adjustbox}
\usepackage{xcolor}
\usepackage{fontspec}
\usepackage[authoryear,round]{natbib}
\usepackage{hyperref}
\hypersetup{
    colorlinks=true,
    linkcolor=blue,
    filecolor=magenta,
    urlcolor=cyan,
}
\usepackage{caption}
% Removes the "Figure:" prefix from captions.
\captionsetup[figure]{name=}

\graphicspath{{../Figures/}{./Figures/}}

% ------------------------------------------------------------------------------
% Palette. Emphasis uses the two muted text colors below. The named colors let
% TikZ diagrams reuse the same palette.
% ------------------------------------------------------------------------------
\definecolor{houseblue}{RGB}{31,58,110}    % matches blue!70!black
\definecolor{housered}{RGB}{140,27,27}     % matches red!70!black

% ------------------------------------------------------------------------------
% Emphasis
%   \term{...}   a defined term or key phrase, bold muted blue
%   \alert{...}  a warning or the one thing to notice, bold muted red
%   \muted{...}  secondary text, gray
%   \code{...}   inline code
% ------------------------------------------------------------------------------
\newcommand{\term}[1]{\textbf{\textcolor{blue!70!black}{#1}}}
\renewcommand{\alert}[1]{\textbf{\textcolor{red!70!black}{#1}}}
\newcommand{\muted}[1]{\textcolor{gray}{#1}}
\newcommand{\code}[1]{\texttt{#1}}

% ------------------------------------------------------------------------------
% Title frame
%   \titleframe{Title}{Authors and affiliations}{Venue and date}{email}
% ------------------------------------------------------------------------------
\newcommand{\titleframe}[4]{%
  \begin{frame}[c]
    \centering
    \huge{#1}\\
    \vspace{0.25in}
    \normalsize{#2}\\
    \vspace{0.25in}
    \footnotesize{#3}\\
    \vspace{.15in}
    \footnotesize{This Version: \today}\\
    \vspace{.25in}
    \href{mailto:#4}{#4}
  \end{frame}%
}

% ------------------------------------------------------------------------------
% Figure frames
% ------------------------------------------------------------------------------

% \source{...}  small gray attribution line under an image
\newcommand{\source}[1]{%
  \par\vspace{2pt}%
  {\scriptsize\textcolor{gray}{Source: #1}\par}%
}

% \figframe{Title}{graphic}{caption}   one figure in a titled frame
\newcommand{\figframe}[3]{%
  \begin{frame}{#1}
    \begin{center}
      \includegraphics[width=\textwidth,height=0.80\textheight,keepaspectratio]{#2}
    \end{center}
    {\scriptsize\textcolor{gray}{#3}\par}
  \end{frame}%
}

% \fullfigframe{graphic}{caption}      one figure, no title, maximum size
\newcommand{\fullfigframe}[2]{%
  \begin{frame}[plain]
    \begin{center}
      \includegraphics[width=\textwidth,height=0.90\textheight,keepaspectratio]{#1}
    \end{center}
    {\scriptsize\textcolor{gray}{#2}\par}
  \end{frame}%
}

% ------------------------------------------------------------------------------
% keybox. One quiet box for the single takeaway of a frame. Use it rarely.
% ------------------------------------------------------------------------------
\newtcolorbox{keybox}{
  colback=black!3, colframe=orange!55!black,
  boxrule=0.6pt, arc=1pt,
  left=6pt, right=6pt, top=4pt, bottom=4pt
}

% ------------------------------------------------------------------------------
% Listings
% ------------------------------------------------------------------------------
\lstset{
  basicstyle=\ttfamily\scriptsize,
  keywordstyle=\color{blue!70!black}\bfseries,
  commentstyle=\color{gray}\itshape,
  stringstyle=\color{green!40!black},
  frame=single,
  framerule=0.4pt,
  rulecolor=\color{gray!50},
  breaklines=true,
  numbers=none,
  showstringspaces=false,
  tabsize=2
}

% ------------------------------------------------------------------------------
% Section dividers. Each \section{...} produces a plain frame with its title.
% ------------------------------------------------------------------------------
\AtBeginSection[]{%
  \begin{frame}[plain,c]
    \begin{center}
      {\usebeamercolor[fg]{structure}\Large\bfseries\insertsection}
    \end{center}
  \end{frame}%
}
